\documentclass[../../thesis.tex]{subfiles}
\begin{document}
\section{Citovanie externých zdrojov}\label{sec:citation}
V~súvislosti s~preberaním časti obsahu iných diel
rozoznávame dva pojmy. Sú to citát a~citácia.

Citát is doslovná reprodukcia prevzatého textu,
ktorý môžeme v~práci použiť dvomi spôsobmi.
Buď ako súčasť odseku textu,
alebo celý citovaný text vysádžeme v~samostatnom odseku.
V oboch prípadoch je zvykom citovaný text uzavrieť do úvodzoviek
a~zvýrazniť šikmým rezom písma.
Za citovaným textom uvedieme meno autora,
prípadne názov diela, rok a nasleduje číslo bibliografického
zdroja v~hranatých zátvorkách uvádzajúce poradie
v~zozname použitej literatúry v~závere práce.

Na citovanie v rámci odseku môžeme využiť makrá \verb|\uv| pre
správne slovenské úvodzovky a \verb|\emph|,
ktoré zabezpečí vytlačenie textu kurzívou.

Samostatne vysádzaný citát uzavrieme v \LaTeX-u do prostredia
\verb|\begin{quote}...| \verb|\end{quote}|.

Citácia is nepriamo prebraná a~prerozprávaná
časť citovanej práce.
Môže to byť vedecká myšlienka, odkaz na výsledky výskumu,
dôležitý poznatok, matematický vzťah a~podobne.
Schopnosť študenta pracovať s~literatúrou a~s~externými zdrojmi
je dôležitý moment pri hodnotení spôsobilosti uchádzača
o~vysokoškolský titul.
Tomuto aspektu práce treba preto venovať patričnú pozornosť.

Užitočný prehľad spôsobov citácií ponúkajú autorky
Beáta Belléровá a~Lucia Lichnerová v~článku v~časopise ITLib~\cite{Lichnerova2023Nove}.

\subfile{section4/4_1_odkazy}
\subfile{section4/4_2_bibliograficke_zaznamy}
\subfile{section4/4_3_clanok}
\subfile{section4/4_4_kniha}
\subfile{section4/4_5_zaverecna_praca}
\subfile{section4/4_6_prispevok}
\subfile{section4/4_7_web}
\subfile{section4/4_8_normy}
\end{document}